\chapter{Mathematical Derivations of LOB Metrics}
\label{ch:mathematical_derivations}

To ensure complete empirical transparency, this appendix mathematically formally derives the core market microstructure metrics utilized in the Data-Flow Matrices (Chapter 19) and the Non-Parametric Models (Chapter 13).

\section{Order Flow Imbalance (OFI)}

Order Flow Imbalance (OFI) is the primary predictive metric for microsecond price movements. It measures the net pressure exerted by market orders against the resting limit orders at the Best Bid ($P_{bid}^{(1)}$) and Best Ask ($P_{ask}^{(1)}$) levels.

Let $\Delta V_{bid}$ denote the change in volume at the best bid, and $\Delta V_{ask}$ denote the change in volume at the best ask over a discrete tick interval $\Delta t$. The OFI is calculated as:

\begin{equation}
    OFI_t = \frac{\Delta V_{bid} - \Delta V_{ask}}{\Delta V_{bid} + \Delta V_{ask}}
\end{equation}

Where $OFI_t \in [-1, 1]$. An $OFI_t$ near $+1$ indicates absolute buy-side pressure (imminent price increase), while $-1$ indicates absolute sell-side pressure (imminent price drop).

\section{Volume-Weighted Average Price (VWAP)}

Because market orders consume liquidity across multiple levels of the Limit Order Book (causing slippage), the execution price of a massive block trade is not the Best Ask or Best Bid, but the Volume-Weighted Average Price.

If a market buy order of size $Q$ consumes liquidity across $N$ price levels ($L_1, L_2 \dots L_N$), with each level providing volume $v_i$ at price $p_i$, the VWAP is given by:

\begin{equation}
    VWAP = \frac{\sum_{i=1}^{N} p_i \cdot v_i}{Q}
\end{equation}

Subject to the constraint that total consumed volume equals the order size:
\begin{equation}
    \sum_{i=1}^{N} v_i = Q
\end{equation}

\section{Bid-Ask Spread and Liquidity Voiding}

The nominal spread $S$ at time $t$ is simply $P_{ask, t}^{(1)} - P_{bid, t}^{(1)}$. 

However, during a flash crash, if the entire Bid side of the order book is depleted (a Liquidity Void), $P_{bid}^{(1)}$ mathematically ceases to exist (or drops to 0). In AMPS, the engine defines a Liquidity Void trigger ($\Lambda$) as:

\begin{equation}
    \Lambda_t = \begin{cases} 
      1 & \text{if } \sum_{i=1}^{d} V_{bid}^{(i)} = 0 \\
      0 & \text{otherwise}
   \end{cases}
\end{equation}

When $\Lambda_t = 1$, the LULD circuit breakers detailed in Chapter 3 are forcefully engaged by the Numba kernel.
